% Generated draft: requires mathematical and source review.
\documentclass[11pt]{article}
\usepackage[margin=25mm]{geometry}
\usepackage{fontspec}
\setmainfont{Latin Modern Roman}
\usepackage{amsmath,amssymb,mathtools,unicode-math}
\setmathfont{Latin Modern Math}
\usepackage{enumitem,needspace,xurl,hyperref}
\setlength{\parindent}{0pt}
\setlength{\parskip}{5pt}
\setlength{\emergencystretch}{3em}
\newcommand{\N}{\mathbb{N}}
\newcommand{\Z}{\mathbb{Z}}
\newcommand{\Q}{\mathbb{Q}}
\newcommand{\R}{\mathbb{R}}
\newcommand{\C}{\mathbb{C}}
\newcommand{\F}{\mathbb{F}}
\newcommand{\A}{\mathbb{A}}
\newcommand{\PP}{\mathbb{P}}
\newcommand{\E}{\mathbb{E}}
\newcommand{\eps}{\varepsilon}
\newcommand{\dd}{\,\mathrm{d}}
\begin{document}
\section*{UnsolvedMath ID: 30000284}
\textbf{Source problem code:} catalogue entry
\subsection{Definitions and mathematical statement}
\textbf{Three-part problem.} For integers \(2\le k<n/2\), characterize compact \(n\)-manifolds admitting a \(k\)-positive metric, identify topological obstructions, and determine when a conformal class containing such a metric also contains a metric of constant \(\sigma_k\)-curvature. The strict upper bound on \(k\) is retained, so the stated range implies \(n\ge5\).\par In the standard Schouten-tensor convention, for a smooth Riemannian metric \(g\) in dimension \(n\ge3\), define\[A_g=\frac1{n-2}\left(\operatorname{Ric}_g-\frac{R_g}{2(n-1)}g\right),\qquad \sigma_k(g)=\sigma_k\bigl(\lambda(g^{-1}A_g)\bigr),\] where \(\sigma_j(\lambda)=\sum_{i_1<\cdots<i_j}\lambda_{i_1}\cdots\lambda_{i_j}\) is the \(j\)th elementary symmetric polynomial. The positive G\aa rding cone is \(\Gamma_k=\{\lambda\in\R^n:\sigma_j(\lambda)>0\text{ for }1\le j\le k\}\), equivalently the component of \(\{\sigma_k>0\}\) containing the positive orthant. A \(k\)-positive, or \(k\)-admissible, metric has \(\lambda(g^{-1}A_g)\in\Gamma_k\) at every point. This standard interpretation must be checked because the source does not itself define \(k\)-positive; positivity of \(\sigma_k\) alone does not specify the elliptic branch.\par Under the usual closed-manifold interpretation, the conformal existence question asks: if \(g\) is \(k\)-admissible, is there a smooth \(u:M\to\R\) and a constant \(c>0\) for which \(\widehat g=e^{2u}g\) remains \(k\)-admissible and \(\sigma_k(\widehat g)\equiv c\)? More generally, determine the precise conditions on the conformal class ensuring such a solution. No volume normalization, prescribed value of \(c\), or local-conformal-flatness assumption is supplied. The source says compact but does not explicitly exclude boundary; boundary problems would need additional boundary data. The separate metric-existence and topological-obstruction tasks are not reduced to this conformal PDE.\par The supplied triage checked 2026-08-21 labels the problem partially solved. It reports broad existence and compactness results in locally conformally flat and low-\(k\) cases, with general non-locally-conformally-flat cases described as unresolved. The exact theorem scopes and any restrictions within the stated \(2\le k<n/2\) range are not reproduced, so no blanket assertion for all such cases is inferred. This dated assessment requires independent verification.
\subsection{Short English statement}
Classify manifolds with Schouten \(k\)-positive metrics, find their topological obstructions, and solve the constant-\(\sigma_k\) problem within admissible conformal classes. The specified range is strictly \(2\le k<n/2\), and the positive elliptic cone must be retained.
\subsection{Sources}
\begin{itemize}\item Local UnsolvedMath registry, permanent ID 30000284, code catalogue entry; \emph{Topology and Constant \(\sigma_k\)-Curvature Metrics}, Oberwolfach Reports Open Problems.\item \emph{Differentialgeometrie im Grossen (2006)}, Oberwolfach Reports, DOI 10.4171/owr/2005/35, pp. 1983--2042. \url{https://doi.org/10.4171/owr/2005/35}. Proposed year and year tag: 2005; no proposer named.\item Supplied literature label \emph{original Oberwolfach report}: \url{https://doi.org/10.4171/owr/2005/35}.\item Supplied literature label \emph{Conformal metrics with prescribed curvature}: \url{https://math.osu.edu/~guan/Papers/guan2007-ajm.pdf}.\item Literature triage checked 2026-08-21, status \texttt{partially\_solved}; record updated 2026-09-16. The \texttt{corrected\_verified} field reports primary-source checking, not independently verified here.\end{itemize}
\end{document}
